% TOP_PROBLEM: {"id": 3055, "title": "Extremal $4$-Neighbour Bootstrap Percolation in the Hypercube", "category_id": 2, "category": {"id": 2, "name": "combinatorics", "display_name": "Combinatorics", "description": "Counting problems, graph theory, discrete structures.", "slug": "combinatorics", "order_index": 2, "created_at": "2026-07-31T15:26:25.671Z"}, "status": "open", "problem_number": "OPG-60003", "set_id": 12}
% FIRST_POSTED: 2026-10-01
% ENTRY_KIND: statement_only
% UnsolvedMath ID 3055; source catalogue code OPG-60003
% !TeX program = lualatex
% Definition and sources only; this file is not an LLM solution attempt.
\documentclass[11pt,a4paper]{article}
\usepackage[margin=25mm]{geometry}
\usepackage{fontspec}
\setmainfont{lmroman10-regular.otf}[BoldFont=lmroman10-bold.otf,
 ItalicFont=lmroman10-italic.otf,BoldItalicFont=lmroman10-bolditalic.otf]
\usepackage{amsmath,amssymb,mathtools}
\usepackage{unicode-math}
\setmathfont{latinmodern-math.otf}
\AtBeginDocument{\let\setminus\smallsetminus}
\usepackage{xurl,hyperref}
\hypersetup{colorlinks=true,urlcolor=blue}
\setcounter{secnumdepth}{0}
\setlength{\parindent}{0pt}
\setlength{\parskip}{5pt}
\setlength{\emergencystretch}{3em}
\begin{document}
\section*{Extremal $4$-Neighbour Bootstrap Percolation in the Hypercube}\label{3055}
{\small UnsolvedMath ID 3055 · OPG-60003 · Combinatorics · OpenGarden}
\subsection{Definitions and mathematical statement}
For an integer $d\ge4$, the hypercube graph $Q_d$ has vertex set
$\{0,1\}^d$; two vertices are adjacent precisely when their binary strings
differ in one coordinate. Write $N(v)$ for the $d$ neighbours of $v$.
Choose an initial infected set $A\subseteq V(Q_d)$, and recursively define
\[
 A_0=A,\qquad
 A_{t+1}=A_t\cup\{v\in V(Q_d):|N(v)\cap A_t|\ge4\}
 \quad(t\ge0).
\]
Infections are permanent, and each step is simultaneous. Since the graph
is finite and the sets increase, the process eventually stabilizes.
The set $A$ percolates when $\bigcup_{t\ge0}A_t=V(Q_d)$.

Define the extremal number
\[
 m(Q_d,4)=\min\{|A|:A\subseteq V(Q_d),\ A\text{ percolates}\}.
\]
The problem is to determine this integer exactly as a function of $d$.
A complete determination needs both a construction of an initial set of
the claimed cardinality for each dimension and a lower bound excluding
all smaller sets. No randomness is imposed on the initial set, and the
percolation time is not the quantity being minimized.

The threshold is exactly four infected neighbours, regardless of dimension;
therefore this is the four-neighbour process rather than a majority rule.
For completeness, when $0\le d<4$ no initially healthy vertex can become
infected, so the same definition gives $m(Q_d,4)=2^d$. The nontrivial range
of the source problem is $d\ge4$.
\subsection{Short English statement}
How few initially infected vertices can eventually infect the entire $d$-dimensional cube when four infected neighbours are needed?
\subsection{Sources}
\begin{itemize}
\item[S1] ULAM.AI and the UnsolvedMath Contributors, supplied problems.json, record 3055 (OPG-60003), OpenGarden collection. Catalogue identity and source transcription; CC BY 4.0.
\url{https://huggingface.co/datasets/ulamai/UnsolvedMath}.
\item[S2] Open Problem Garden, Extremal $4$-Neighbour Bootstrap Percolation in the Hypercube. Original problem and discussion; the mathematical formulation below is expanded from the supplied source record.
\url{http://www.openproblemgarden.org/op/extremal_4_neighbour_bootstrap_percolation_in_the_hypercube}.
\end{itemize}
\end{document}
